\section{Subsidies in an Open Economy}

\begin{frame}{Subsidies I - Subsidising Production}
    \showmytikz[1]{../shared/tikz/tikz-subsidy-trade.tex}
    \mycomment{
        Subsidies are financial assistance to producers by the government, lowering their cost of production and engaging higher output. There are two types: production subsidies and export subsidies. Here, we look at production subsidies. These lower the cost of the producer for both the world market and the internal market. The aim will be to reduce imports.
    }
\end{frame}
\begin{frame}{Subsidies I - Subsidising Production}
    \showmytikz[1,2]{../shared/tikz/tikz-subsidy-trade.tex}
    \mycomment{
    The subsidy has shifted the supply curve to the right. However, \textit{real} costs remain the same; only a part of it is now covered by the government. 

    Consumer surplus remains the same, as price and quantity are left unchanged. They will, however, now consume more local produce.

    This subsidy may reduce producer's incentives to become more cost efficient: the state shoulders the extra costs. Since the costs are not going down, producer surplus is not increased.
    }
\end{frame}

\begin{frame}{Subsidies I - Subsidising Production}
    \showmytikz[1,2,3]{../shared/tikz/tikz-subsidy-trade.tex}
    \mycomment{
        However, the producer surplus has increased due to the subsidy increasing domestic prices. 
    }
\end{frame}

\begin{frame}{Subsidies I - Subsidising Production}
    \showmytikz[1,2,3,4]{../shared/tikz/tikz-subsidy-trade.tex}
    \mycomment{
        Of course, the subsidy must have come from somewhere. If it is financed by taxes, the loss of (consumer) income will (in some markets) shift the demand curve (to the left), reducing consumer surplus. There is also an excess burden due to the costs of tax collection and subsidy payout.
        
        There is also a welfare loss: that part of the subsidy which exceeds the gain in producer surplus (dead weight loss, DWL). 
    }
\end{frame}


\begin{frame}{Subsidies II - Export Subsidy}
    \showmytikz[1]{../shared/tikz/tikz-exportsubsidy.tex}
    \mycomment{
        A country exports (already), because there are more people outside the country willing to pay the price for the goods than inside the country.    

        The government subsidises only exports: the supply curve remains the same, but for every exported unit local producers receive an extra payment.

        If the country is small, this does not affect the world price.
    }
\end{frame}

\begin{frame}{Subsidies II - Export Subsidy}
    \showmytikz[1,2]{../shared/tikz/tikz-exportsubsidy.tex}
    \mycomment{By supporting the export, producers will produce more at the higher price. The consumer inside the country will now also have to pay the higher price.}
\end{frame}

\begin{frame}{Subsidies II - Export Subsidy}
    \showmytikz[1,2,3]{../shared/tikz/tikz-exportsubsidy.tex}
    \mycomment{The size of the subsidy is larger than the `gain' of welfare, so there is a dead weight loss. 
    
    On top, of course, the government must raise the funds for the subsidy elsewhere, by taxing either another industry or through income taxes. This has of course detrimental effects on the other industries and to internal demand, which is not modelled here.

    }
\end{frame}

\begin{frame}{Subsidies II - Export Subsidy}
    \mycomment{
    So far, we have looked at a small country, so that the subsidy does not affect the world price. If, however, the quantity the country (now) exports is large, global prices might indeed fall. 
    
    As a result, consumers in the rest of the world profit from lower prices. The export subsidy then effectively "sponsors" welfare gains of consumers abroad, and results in welfare loss of producers abroad.
    }
    \showmytikz[1,2,3,4]{../shared/tikz/tikz-exportsubsidy.tex}
\end{frame}


\begin{frame}{Subsidies III - From Import to Export by Subsidy}
    \showmytikz[1]{../shared/tikz/tikz-exportsubsidy2.tex}
    \mycomment{Here, the original difference between the world price and the local cost structure is so that the country is an importer. However, the government subsidies so much that the situation  inverses - clearly a costly and inefficient affair. 
    
    }
\end{frame}

\begin{frame}{Subsidies III - From Import to Export by Subsidy}
    \showmytikz[1,2]{../shared/tikz/tikz-exportsubsidy2.tex}
    \mycomment{
    The subsidy is pushing a low domestic price so much up that the country becomes an exporter, even though it should be an importer.

    }
\end{frame}

\begin{frame}{Subsidies III - From Import to Export by Subsidy}
    \showmytikz[1,2,3]{../shared/tikz/tikz-exportsubsidy2.tex}
    \mycomment{
    A large government payout increases local producers' surplus at the costs of local consumers, but there is a large dead weight loss: the subsidy is much larger than the part which `only' shifts surplus from consumers to producers. 

    This all is under the assumption that this aggressive policy does not change the world market price.
    }
\end{frame}

\begin{frame}{Non-tariff barriers}
    \mytext{
    \begin{itemize}
        \item requirements
        \item restrictions (weapons etc.)
        \item customs delays
        \item technical barriers
        \item licensing, standards
        \item packaging, labeling, localisation
        \item sanitation rules
    \end{itemize}
    }
    \mycomment{Non-tariff barriers typically affect internal and external suppliers - otherwise, they would be unfair practice.    
    }
\end{frame}

\begin{frame}{Public Procurement}
    \mytext{
    \begin{itemize}
        \item government to buy only locally produced goods 
        \item forbidden under WTO rules
    \end{itemize}
    }
    
\end{frame}

\begin{frame}{Embargoes and Sanctions}
    \mytext{
    Außenwirtschaftsgesetzt AWG: trade is free (in principle)
    \begin{itemize}
        \item trade barriers
        \item asset freezes
        \item travel bans
        \item arms embargos
        \item restriction on financial transactions
        \item international sanctions vs. unilateral coercive measures
        \item blockades (acts of war)
        \item financial system restrictions (SWIFT etc.)
        \item EU Instex (barter / clearing), \href{https://eur-lex.europa.eu/eli/reg/1996/2271/oj/eng}{EU Blocking Regulation}
        \item US: OFAC (Office of Foreign Asset Control)
    \end{itemize}
    }
    \mycomment{
    Embargoes and sanctions restrict trade in respect to
    \begin{itemize}
        \item countries
        \item specific goods
        \item persons or organisations.
    \end{itemize}
    
    This typically shifts trade patterns towards third / neutral countries.

    Unilateral sanctions, however, tend to reduce GDP growth of targeted countries \cite{neuenkirchneumeier}.
    
    For an overview see \cite{gsdb}.
   
    }
\end{frame}

\begin{frame}{Retaliation}
    \mytext{
    \begin{itemize}
        \item trade efficient: tariffs inefficient
        \item large countries vs.\ small country
        \item optimum tariff?
        \item pressure groups, populism...
        \item with or without retaliation?
    \end{itemize}
    }
\end{frame}

\begin{frame}{Cooperation or Defection?}
    \mytext{
    \begin{itemize}
        \item cooperation would be best for all
        \item but: defect from cooperation yields (short-term) benefit
        \item cooperation not stable solution
        \item Game Theory
    \end{itemize}
    }
\end{frame}


\begin{frame}{Prisoner's Dilemma}
    
    \showmytikz{../shared/tikz/tikz-prisonersdilemma-trade.tex}
\end{frame}

\begin{frame}{Borders and Economic Areas}
    \mytext{
    \begin{itemize}
        \item free trade is always welfare maximising
        \item trade barreers / tariffs most effective where there is a lot of trade (bridges, ports)
        \item actual trade flows and national borders need not correspond
        \item e.g. North America: borders East-West, trade North-South
        \item borders often militarily / historically given
        \item optimum currency area is a different question
        \item traders (in ports, global cities) often international
        \item policy often national / rural / irrational
    \end{itemize}
    }
\end{frame}

\begin{frame}{National Trade Policy}
    \mytext{
    \begin{itemize}
        \item trade is generally, globally, good
        \item specialisation is good
        \item winners and losers to trade
        \item consumers less organised than producers
        \item produce for exports or for local markets?
        \item traders are also a pressure group
        \item specialisation has also disadvantages
    \end{itemize}
    }
\end{frame}

